\section{HPO and idea-sprint verdicts --- cheap, trustworthy negatives}

\subsection{What this note covers}

Four experiments that all said ``no''. Or ``no measurable change''.
LGBM hyperparameter tuning. A rolling spike threshold.
Two conformal finite-sample fixes.
None moved the headline MAE. All are worth writing down.
The meta-lesson at the end is the real point of the note.

\subsection{LGBM HPO --- defaults survive}

The champion ran ``conservative untuned defaults'' since M4.
Owner authorized a proper tuning campaign.
14 configs, 1095-day windows. Screen on test year 1.
Confirm the top 3 plus control on test year 2.
The confirm year is never used for selection --- no peeking.
Pre-declared gate: beat control by $\ge 0.10$ MAE on \emph{both} years.

\begin{tabular}{lccl}
\hline
Config & Screen MAE & Confirm MAE & Verdict \\
\hline
control (defaults) & 17.62 & 16.98 & stays \\
lr02\_n1500 (best) & 17.55 & 16.91 & gate failed \\
lr03\_n1000 & 17.54 & 16.97 & gate failed \\
leaves127 & 17.62 & 17.10 & gate failed \\
\hline
\end{tabular}

\vspace{0.5em}
No config clears the gate. Best gain is 0.02--0.08 MAE per year.
The defaults were already near the optimum for this feature set.
The HPO surface is flat. That is the normal LGBM outcome
on well-engineered features. Chasing 0.05 MAE buys fragility,
not skill. The model card now reads
``tuned 2026-07-25 --- defaults survived''.

\subsection{Side finding --- a third sign flip}

Dropping the load-lag features helped at 365-day windows:
$-0.12$ MAE. At 1095-day windows it \emph{hurts}: $+0.09$
(17.71 vs 17.62). Same feature, opposite verdict, different window.
This is the third documented ablation sign flip in the project.
Lesson: an ablation verdict is only valid for its training window.
Always quote the window with the verdict.

\subsection{Rolling spike threshold --- rejected}

The spike classifier flags the top-5\% price hours.
It uses a static threshold from the training window.
Idea: make the threshold roll over a trailing 90 days instead.
Result: AUC 0.955 vs 0.966. Precision@2 down too.
The static train-window threshold stays.
The rolling version reacts to noise, not to regime.

\subsection{Conformal finite-sample fixes --- adopted, effect $\sim$0}

The validation review flagged two statistical defects (E3, E4).
The conformal quantile was linearly interpolated.
The asymmetric path sized the upper tail from the lower-tail count.
Both can slightly under-deliver the coverage guarantee.

Fix: use \texttt{np.quantile(..., method="higher")} and size each
tail from its own score count. We adopted both in \texttt{src}.
Measured effect on our numbers: about zero.
At our window size (2{,}160 scores) interpolation barely matters.
We adopted them anyway --- for the finite-sample \emph{guarantee},
not for a MAE gain. A correctness fix does not need to move a metric.

\subsection{The meta-lesson --- pre-declared gates make negatives cheap}

Every experiment here declared its gate \emph{before} running.
HPO: beat control by $\ge 0.10$ on both years.
Spike threshold: beat AUC 0.966. Conformal: restore the guarantee.

Because the bar was set first, a ``no'' is a clean result.
No temptation to move the goalposts after seeing the numbers.
No cherry-picking the one year or one metric that looks good.
A pre-declared gate turns a negative into evidence you can trust.
It also makes negatives cheap to report: the gate is unmet, done.
Most of a real desk's experiments fail. This is how you fail honestly.

\subsection{Interview line}

``I ran a 14-config LGBM tuning campaign with a screen year and a
held-out confirm year, gate declared up front: beat the defaults by
0.10 MAE on both. Nothing cleared it --- the surface was flat, which
is the honest LGBM outcome on engineered features. Same discipline
killed a rolling spike threshold and confirmed two conformal fixes
had zero measured effect. Pre-declaring the gate is what makes a
negative result cheap to run and safe to trust.''
